% Modernised reading — fonds Alexandre Grothendieck, shelfmark 116, pages 1-10.
%
% This is an INTERPRETATION, not a transcription. Everything the critical
% apparatus carried moves into footnotes. In any dispute of substance the
% transcription governs: whoever cites Grothendieck must cite batch-01.fr.tex.
%
% Pass: Opus 5 (claude-opus-5), 2026-09-13 — first pass, unchecked against the
% pages by a human. The folder was read whole; it holds one batch, pages 1-10.
% Page 1 is a wrapper (« Mapping cone »), page 8 blank, page 10 an
% administrative notice of 19 June 1981; they carry no mathematics.
%
% What the folder is. Under his title « Relations entre morphismes de topos et
% recollements de topos » (pages 3-7): a geometric morphism f : X_0 -> X_1 is
% the same thing as a topos X with an open subtopos X_0 and closed complement
% X_1 glued along a functor with a left exact left adjoint. X = C(f), which
% he calls the mapping-cone, is the Artin gluing along f_*; its closed part is
% a retract with the same cohomology; the relative cohomology of the pair is
% the cohomology « of X_1 mod X_0 »; and for presheaf topoi the gluing datum
% is a profunctor. Page 2 (the wrapper's back) and page 9 are separate notes.
%
% Notation fixed for the whole document. Script X, X_0, X_1 topoi, F(X) the
% category of sheaves (as he writes it); roman X, X_0, X_1 small categories
% (pages 6-7) and X^ their presheaf topoi Top(X). Sheaves on the gluing are
% triples (F_0, F_1, u), u : f^*F_1 -> F_0. j : X_0 -> X the open inclusion,
% i_1 : X_1 -> X the closed one, r : X -> X_1 the retraction (the reading's
% name; the leaf's margin describes its inverse image and calls it i_*).
% Top(Delta_1) the Sierpinski topos.
%
% Substantive departures from the pages, each footnoted where it occurs:
% — page 6: a functor commuting with all limits has a LEFT adjoint; the leaf
%   says « à droite » there and « à gauche » on pages 4 and 7.
% — page 3: « mapping-cone »: the topos is the mapping cylinder (his own
%   margin writes it (X_0 x Delta_1^) ⊔_{X_0} X_1); its relative cohomology
%   mod X_0 is the cohomology of the cone. The reading keeps his word and says
%   so.
% — page 3, the margin's « i_*(F_1) = (f^*(F_1), F_1, id) » is the inverse
%   image of the retraction r, left adjoint to i_1^*, not a direct image.
% — page 5: the N.B. writes (F_1, f^*F_1, u) in the reverse of the order fixed
%   on page 3.
% — page 9: the isomorphism of limits needs a hypothesis on A -> B that the
%   leaf does not state; the reading gives one under which it holds.
%
% Modern names given and footnoted as ours: Artin gluing (recollement, SGA 4
% IV 9.5), mapping cylinder of topoi, Sierpinski topos classifying open
% subtopoi, local cohomology / cohomology with supports, collage (cograph)
% of a profunctor, weak homotopy equivalence of topoi, Quillen's Theorem A
% type cofinality for fibres versus comma categories.
%
% What the folder announces and never establishes: the converse asked on
% page 4 (« réciproque ? »); the conditions on (X_0, X_1, phi) making both
% inclusions weak equivalences, which page 6 sets out to find and does not.
\documentclass[11pt,a4paper]{article}
% Le préambule commun à toutes les transcriptions du fonds Grothendieck.
%
% Il tient en une idée : **l'incertitude se compose, elle ne se résout pas.**
% Une transcription qui lisse un mot douteux a détruit la seule chose pour
% laquelle on la faisait — un lecteur doit pouvoir savoir, à chaque endroit,
% ce qui était sur le feuillet et ce qui a été ajouté. Les six macros qui
% suivent sont donc l'appareil critique, et non de la décoration.
%
% Les mêmes commandes sont reconnues par `scripts/render.mjs`, qui produit la
% vue de lecture du site. Une macro ajoutée ici doit l'être aussi là-bas,
% faute de quoi le rendu échouera bruyamment — ce qui est voulu : un rendu
% silencieusement faux est pire qu'un rendu absent.

\ProvidesPackage{grothendieck}[2026/08/08 Transcriptions du fonds Grothendieck]

\usepackage{amsmath,amssymb,amsthm}
\usepackage{mathrsfs}
\usepackage{stmaryrd}
% Les diagrammes commutatifs sont partout dans ce fonds : c'est la langue
% même de la géométrie algébrique de Grothendieck, et une transcription qui
% les rendrait en prose ne transcrirait rien.
\usepackage{tikz-cd}
% Français par défaut — c'est la langue des feuillets — mais l'anglais est
% chargé aussi : la traduction et le résumé sont des documents anglais, et la
% typographie française (espaces avant les deux-points) y serait une faute.
% Ces documents basculent par \selectlanguage{english} en tête de corps.
\usepackage[english,french]{babel}
\usepackage{fontspec}
\usepackage{geometry}
\geometry{a4paper,margin=25mm,marginparwidth=28mm}
\usepackage{marginnote}
\usepackage{xcolor}
% ulem, not soul. `soul`'s \st and \ul are built on a character-by-character
% scan that XeLaTeX's font machinery breaks: the document fails with an
% undefined control sequence rather than degrading. ulem does the same two jobs
% and is Unicode-engine safe. `normalem` stops it hijacking \emph, which the
% transcriptions use for Grothendieck's own emphasis.
\usepackage[normalem]{ulem}
\usepackage{hyperref}
\hypersetup{colorlinks=true,linkcolor=black,urlcolor=black,pdfborder={0 0 0}}

% --- Métadonnées du lot -----------------------------------------------------
% Reprises telles quelles par le script de rendu, qui les affiche en tête de
% la vue de lecture. Elles fixent ce que le fichier prétend couvrir : sans
% elles, une transcription flotte sans point d'attache dans un fonds de seize
% mille feuillets.

\newcommand{\folder}[1]{\gdef\grothFolder{#1}}
\newcommand{\batch}[1]{\gdef\grothBatch{#1}}
\newcommand{\leaves}[2]{\gdef\grothFirst{#1}\gdef\grothLast{#2}}
\newcommand{\foldertitle}[1]{\gdef\grothFoldertitle{#1}}
\gdef\grothFolder{?}\gdef\grothBatch{?}\gdef\grothFirst{?}\gdef\grothLast{?}\gdef\grothFoldertitle{}

% --- Le résumé --------------------------------------------------------------
%
% Il ouvre la lecture modernisée, sous le titre « Résumé », et il est composé
% en petit. Ce n'est pas un ornement : c'est le seul passage du document qui
% ne soit pas des mathématiques, et un lecteur doit pouvoir le reconnaître
% avant de l'avoir lu — puis le sauter s'il connaît déjà le sujet, ou s'y
% arrêter s'il ne le connaît pas. Le filet qui le suit marque la reprise du
% corps.

\newenvironment{resume}
  {\small\setlength{\parskip}{0.45em}}
  {\par\medskip\hrule\medskip}

% --- Mots-clés ---------------------------------------------------------------
%
% En anglais, en clôture du résumé de la lecture modernisée : le vocabulaire
% moderne sous lequel chercher ce que les feuillets construisent. Le manifeste
% du site (`npm run manifest`) les extrait pour étiqueter la cote entière —
% cette ligne est la source unique des tags, il n'y a pas de fichier à côté.
\newcommand{\keywords}[1]{\par\smallskip\noindent
  {\footnotesize\textsc{Keywords} --- \textit{#1}}\par}

% --- L'appareil critique ----------------------------------------------------

% Le numéro de feuillet, en marge. C'est l'ancre qui permet de régler un
% désaccord sur une formule en regardant *un* feuillet plutôt qu'une chemise
% de six cents. Le site s'en sert aussi pour tourner les pages du fac-similé
% quand on fait défiler la transcription.
\newcommand{\leaf}[1]{%
  \marginnote{\footnotesize\color{gray!70}\textsf{#1}}%
  \label{leaf:#1}\ignorespaces}

% Illisible : on ne devine pas. Un mot inventé qui se lit comme les autres est
% le pire résultat possible.
\newcommand{\ill}{\textcolor{red!60!black}{[\,\ldots\,]}}

% Lecture douteuse : le mot est donné, et signalé comme incertain.
\newcommand{\uncertain}[1]{\uline{#1}}

% Ajout de l'éditeur — une lettre manquante, un mot sous-entendu.
\newcommand{\add}[1]{\textcolor{blue!50!black}{[#1]}}

% Biffé par Grothendieck lui-même. Ce qu'il a rayé fait partie du document :
% une rature dit souvent où la pensée a changé de direction.
\newcommand{\struck}[1]{\sout{#1}}

% Note du transcripteur, sur l'état matériel du feuillet.
\newcommand{\note}[1]{\par\smallskip{\small\color{gray!60!black}\textit{[#1]}}\par\smallskip}

% Note marginale de Grothendieck, distincte du corps du texte.
\newcommand{\marginal}[1]{\par\smallskip\noindent\hspace{1em}%
  {\small\color{gray!60!black}#1}\par\smallskip}

% --- Titre ------------------------------------------------------------------

\AtBeginDocument{%
  \begin{center}
    {\large\bfseries Fonds Alexandre Grothendieck --- cote n\textsuperscript{o}~\grothFolder}\\[2pt]
    {\itshape\grothFoldertitle}\\[4pt]
    {\small feuillets \grothFirst--\grothLast\ (lot \grothBatch)}\\[2pt]
    {\footnotesize\color{gray!60!black}Transcription établie d'après le fac-similé numérisé
     par l'Université de Montpellier.\\
     Les lectures incertaines sont soulignées, les illisibles notées [\,\ldots\,],
     les ajouts entre crochets.}
  \end{center}
  \vspace{1em}}

\endinput


\folder{116}
\pages{1}{10}
\dating{[à partir de 1981]}
\watermark{Édition de démonstration\\interprétation personnelle de l'œuvre}
\foldertitle{Mapping-cône etc. : notes manuscrites (s.d.) — lecture modernisée du dossier entier}

\begin{document}

\section*{Résumé}

\begin{resume}
En topologie, une application continue $f : X_{0} \to X_{1}$ se remplace par
une inclusion sans rien perdre de son type d'homotopie : on fabrique le
\emph{cylindre} de $f$, en posant $X_{0} \times [0,1]$ à côté de $X_{1}$ et en
collant le bout $X_{0} \times \{1\}$ sur $X_{1}$ par $f$. Le cylindre se rétracte
sur $X_{1}$, et $X_{0}$ y est plongé à l'autre bout. L'application est devenue
une inclusion, et la cohomologie relative de cette inclusion est celle du
\emph{cône} de $f$.

Les topos --- les « espaces généralisés » de Grothendieck, définis par leurs
faisceaux plutôt que par leurs points --- admettent la même construction, et
ces feuillets l'écrivent. Un topos qui contient un ouvert et le fermé
complémentaire se reconstitue à partir des deux morceaux et d'un foncteur qui
dit comment les faisceaux de l'ouvert « débordent » sur le fermé : c'est le
\emph{recollement}. Le résultat principal des pages 3 et 4 est qu'un morphisme
de topos est la même chose qu'un recollement d'un type particulier --- celui où
le foncteur de recollement a un adjoint à gauche exact à gauche. Le topos ainsi
recollé est le cylindre du morphisme ; son fermé en est un rétracte, avec la
même cohomologie, et la suite exacte de la paire (ouvert, tout) devient la
cohomologie relative du morphisme.

Les pages 5 à 7 prennent le point de vue inverse --- partir d'un topos muni
d'un ouvert, ou d'un recollement quelconque --- et se demandent quand les deux
inclusions sont des équivalences au sens de l'homotopie. Elles finissent sur
le cas le plus concret : quand tout est un topos de préfaisceaux sur une petite
catégorie, se donner le recollement revient à se donner une famille
d'ensembles $h(x_{1}, x_{0})$ qui joue le rôle des flèches allant d'un morceau
à l'autre.

Le vocabulaire est celui du début des années 1980 : il écrit d'abord
« asphérique », puis le barre à deux reprises pour « équivalence faible ».

\keywords{Artin gluing, recollement of topoi, mapping cylinder, Sierpinski topos, open subtopos, cohomology with supports, profunctor collage, weak homotopy equivalence}
\end{resume}

\section*{Le fil du dossier, et les conventions}

\begin{enumerate}
\item[1.] \emph{Au dos de la chemise} (page 2) : l'adjonction triple d'un
foncteur dont l'image inverse est pleinement fidèle.
\item[2.] \emph{Le cylindre d'un morphisme de topos} (pages 3 et 4) :
faisceaux, rétraction, cohomologie, caractérisation.
\item[3.] \emph{Point de vue inverse} (pages 5 à 7) : ouverts et
$\mathrm{Top}(\Delta_{1})$, recollements quelconques, cas des préfaisceaux.
\item[4.] \emph{Un feuillet séparé} (page 9).
\end{enumerate}

\emph{Conventions.} $\mathcal{X}$, $\mathcal{X}_{0}$, $\mathcal{X}_{1}$ sont des
topos, $\mathcal{F}(\mathcal{X})$ la catégorie des faisceaux sur $\mathcal{X}$.
Aux pages 6 et 7, les lettres droites $X$, $X_{0}$, $X_{1}$ sont de petites
catégories et $\widehat{X} = \mathrm{Top}(X)$ le topos des préfaisceaux sur
$X$.\footnote{La distinction typographique est de la transcription ; le
feuillet ne sépare les deux tracés que par la forme de la lettre.} Un
morphisme de topos $f$ a une image directe $f_{*}$ et une image inverse
$f^{*}$, exacte à gauche et adjointe à gauche de $f_{*}$.
$\mathrm{Top}(\Delta_{1})$ est le topos de Sierpiński, des préfaisceaux sur la
catégorie $0 \to 1$.

\emph{Ce que le dossier annonce et n'établit pas} : la réciproque demandée à la
page 4 ; les conditions sur un recollement qui rendent les deux inclusions
équivalences faibles, que la page 6 se propose de chercher.

\pagerange{2}{2}
\section*{Au dos de la chemise (page 2)}

La chemise porte seule, de sa main, « Mapping cone » ; son verso est un
pense-bête. On y lit l'adjonction triple $i_{!} \dashv i^{*} \dashv i_{*}$
associée à un foncteur $i : A_{0} \to A$ entre petites catégories, entre les
topos de préfaisceaux $\widehat{A}_{0}$ et $\widehat{A}$, dans le cas où $i^{*}$
est pleinement fidèle, $i_{*}$ et $i_{!}$ étant alors des localisations. C'est
exactement la condition pour que
\[
i_{!}\, i^{*}(F) \xrightarrow{\;\sim\;} F,
\qquad
F \xrightarrow{\;\sim\;} i_{*}\, i^{*}(F),
\]
la coünité et l'unité des deux adjonctions étant des isomorphismes pour tout
$F$. Suit une flèche qui « induit iso sur $\Gamma$, puis $H^{*}$ ».\footnote{Les
membres de droite de $i_{!} =$ et $i^{*} =$ sont un même glyphe, coupé par le
bord et non lu. Le feuillet porte aussi une rétraction $p_{A} : A \to A_{0}$,
un triangle de flèches $p$, une flèche $p_{A!}\, i_{!}(\cdot) \to p_{A!}(F)$
dont l'argument n'est pas lu, et un « $\mathrm{Ex}_{\alpha}$ » entre deux mots
illisibles. L'abréviation « pl.\ fid. » de la première ligne n'est pas
assurée ; c'est elle qui rend les deux isomorphismes vrais pour tout $F$, et on
la retient pour cette raison. Les exposants $0$ sont écrits sur le $i$.}

\pagerange{3}{4}
\section*{Le cylindre d'un morphisme de topos (pages 3 et 4)}

\emph{La construction.} Soit $f : \mathcal{X}_{0} \to \mathcal{X}_{1}$ un
morphisme de topos. On le regarde comme un topos fibré sur $\Delta_{1}$, de
fibres $\mathcal{X}_{0}$ et $\mathcal{X}_{1}$, et l'on considère le topos total
$\mathcal{X} = C(f)$, dont les faisceaux sont les triples
\[
(F_{0}, F_{1}, u), \qquad F_{0} \in \mathcal{F}(\mathcal{X}_{0}),\
F_{1} \in \mathcal{F}(\mathcal{X}_{1}), \qquad
u : f^{*}(F_{1}) \to F_{0},
\]
ou, ce qui revient au même par adjonction, $F_{1} \to f_{*}(F_{0})$.

Dans $\mathcal{X}$, $\mathcal{X}_{0}$ est un sous-topos ouvert, d'inclusion
$j : \mathcal{X}_{0} \to \mathcal{X}$, $j^{*}(F_{0}, F_{1}, u) = F_{0}$, et
$\mathcal{X}_{1}$ le fermé complémentaire, d'inclusion
$i_{1} : \mathcal{X}_{1} \to \mathcal{X}$, $i_{1}^{*}(F_{0}, F_{1}, u) = F_{1}$.
Le foncteur de recollement du couple (ouvert, fermé), $i_{1}^{*}\, j_{*}$, n'est
autre que
\[
\varphi = f_{*} : \mathcal{F}(\mathcal{X}_{0}) \to \mathcal{F}(\mathcal{X}_{1}),
\]
puisque $j_{*}(F_{0}) = (F_{0}, f_{*}F_{0}, \mathrm{id})$.\footnote{C'est le
recollement d'Artin des topos (SGA~4, IV, 9.5) le long du foncteur exact à
gauche $f_{*}$ ; le nom n'est pas sur la page. Il appelle $\mathcal{X}$ le
« mapping-cone ». Le topos obtenu est plutôt l'analogue du \emph{cylindre} : sa
propre note en marge l'écrit
$\mathcal{X} = (\mathcal{X}_{0} \times \widehat{\Delta_{1}})
\amalg_{\mathcal{X}_{0}} \mathcal{X}_{1}$, $\mathcal{X}_{0} \times
\widehat{\Delta_{1}}$ étant recollé à $\mathcal{X}_{1}$ par $f$ le long de
$\mathcal{X}_{0} \times \{1\}$ ; et c'est la cohomologie relative de
$\mathcal{X}$ modulo $\mathcal{X}_{0}$, ci-dessous, qui est celle du cône. On
garde son mot pour la section, avec cette précision.}

\emph{La rétraction.} Le foncteur
\[
r^{*} : \mathcal{F}(\mathcal{X}_{1}) \to \mathcal{F}(\mathcal{X}),
\qquad G \longmapsto \bigl(f^{*}G,\ G,\ \mathrm{id}\bigr),
\]
est exact à gauche et adjoint à gauche de $i_{1}^{*}$ : un morphisme
$(f^{*}G, G, \mathrm{id}) \to (F_{0}, F_{1}, u)$ est déterminé par sa composante
$G \to F_{1}$. Il définit donc un morphisme de topos
$r : \mathcal{X} \to \mathcal{X}_{1}$, avec $r_{*} = i_{1}^{*}$, $r\, i_{1} =
\mathrm{id}$ et $r\, j = f$.\footnote{La note marginale de la page 3 parle d'une
« rétraction » de $\mathcal{X}$ sur $\mathcal{X}_{1}$ et écrit
« $i_{*}(F_{1}) = (f^{*}(F_{1}), F_{1}, \mathrm{id})$ ». Le foncteur ainsi défini
est l'image inverse de la rétraction, adjoint à gauche de $i_{1}^{*}$, et non
une image directe ; le nom $r$ est le nôtre.}

\emph{Cohomologie.} L'inclusion du fermé est « strictement une équivalence
faible », et ce qu'il écrit après « i.e. » est que, pour tout faisceau abélien $F = (F_{0}, F_{1}, u)$ sur
$\mathcal{X}$,
\[
H^{*}(\mathcal{X}, F) \xrightarrow{\;\sim\;} H^{*}\bigl(\mathcal{X}_{1},
i_{1}^{*}F\bigr) = H^{*}(\mathcal{X}_{1}, F_{1}).
\]
En effet, une section globale de $F$ est un couple $(s_{0}, s_{1})$ avec
$s_{0} = u(s_{1})$, donc $\Gamma(\mathcal{X}, F) = \Gamma(\mathcal{X}_{1}, F_{1})
= \Gamma(\mathcal{X}_{1}, r_{*}F)$ ; et $r_{*} = i_{1}^{*}$ est exact et, étant
adjoint à droite d'un foncteur exact, préserve les injectifs.\footnote{La page
écrit « strictement asphérique », barre « asphérique » et écrit au-dessus
« équiv.\ f. », puis donne l'isomorphisme de cohomologie après un « i.e. ». La
démonstration est la nôtre ; elle établit l'énoncé cohomologique, qui est ce
que la page écrit. L'énoncé homotopique demande en outre une homotopie entre
$i_{1} r$ et l'identité, que le cylindre fournit et que la page ne détaille
pas.}

\textbf{Caractérisation} (page 4). Se donner un morphisme de topos
$\mathcal{X}_{0} \to \mathcal{X}_{1}$ revient au même que se donner un topos
$\mathcal{X}$, un sous-topos ouvert $\mathcal{X}_{0}$ et son fermé complémentaire
$\mathcal{X}_{1}$, de telle façon que le foncteur de recollement $\varphi$, qui
est toujours exact à gauche, commute aux limites projectives quelconques ---
c'est-à-dire admette un adjoint à gauche $f^{*}$ --- et que cet adjoint soit
exact à gauche.\footnote{L'équivalence « commute aux limites quelconques,
i.e.\ admet un adjoint à gauche » vaut entre topos de Grothendieck par le
théorème de l'adjoint spécial, un topos étant complet, bien normé et muni d'une
petite famille cogénératrice. La page demande « réciproque ? » à propos de
l'asphéricité du fermé, avec un point d'interrogation en marge.}

\emph{La suite de la paire.} La suite exacte de cohomologie à supports dans le
fermé,
\[
\cdots \to H^{i}_{\mathcal{X}_{1}}(\mathcal{X}, F) \to H^{i}(\mathcal{X}, F)
\to H^{i}(\mathcal{X}_{0}, F_{0}) \to H^{i+1}_{\mathcal{X}_{1}}(\mathcal{X}, F)
\to \cdots,
\]
s'écrit donc
\[
\cdots \to H^{i}_{\mathcal{X}_{1}}(\mathcal{X}, F) \to H^{i}(\mathcal{X}_{1},
F_{1}) \to H^{i}(\mathcal{X}_{0}, F_{0}) \to H^{i+1}_{\mathcal{X}_{1}}(\mathcal{X},
F) \to \cdots,
\]
et l'on pose
\[
H^{i}\bigl(\mathcal{X}_{1} \bmod \mathcal{X}_{0},\ (F_{0}, F_{1}, u)\bigr)
\overset{\text{déf}}{=} H^{i}_{\mathcal{X}_{1}}(\mathcal{X}, F),
\]
$f$ étant sous-entendu. Pour $F = r^{*}(F_{1})$, c'est-à-dire $F_{0} =
f^{*}F_{1}$ et $u$ tautologique, on note
$H^{i}(\{\mathcal{X}_{1} \bmod \mathcal{X}_{0}\}, F_{1})$ : c'est la cohomologie
relative du morphisme $f$ à coefficients dans $F_{1}$, et la suite ci-dessus
devient la suite longue
\[
\cdots \to H^{i}(\{\mathcal{X}_{1} \bmod \mathcal{X}_{0}\}, F_{1}) \to
H^{i}(\mathcal{X}_{1}, F_{1}) \xrightarrow{f^{*}} H^{i}(\mathcal{X}_{0},
f^{*}F_{1}) \to \cdots.
\]

\pagerange{5}{7}
\section*{Le point de vue inverse (pages 5 à 7)}

\emph{Équivalences faibles.} Pour que $f$ soit une équivalence faible, il faut
et il suffit que l'immersion ouverte $j : \mathcal{X}_{0} \to \mathcal{X}$ le
soit : on a $f = r\, j$, et $r$ est une équivalence faible dès que $i_{1}$ en
est une, puisque $r\, i_{1} = \mathrm{id}$.\footnote{La page écrit l'énoncé,
barrant encore « asphérique » pour « équiv.\ faible » ; la raison par
$f = rj$ est la nôtre, et elle repose sur l'énoncé homotopique de la page 3,
qui n'est vérifié ici qu'en cohomologie.}

\textbf{N.B.} Les faisceaux localement constants sur $\mathcal{X}$ sont ceux de
la forme $(f^{*}F_{1}, F_{1}, u)$ avec $u$ tautologique et $F_{1}$ localement
constant sur $\mathcal{X}_{1}$ --- autrement dit, ceux de la forme $r^{*}F_{1}$.
C'est ce qui rend $r$ compatible aux systèmes locaux, et pas seulement à la
cohomologie.\footnote{La page écrit $(F_{1}, f^{*}(F_{1}), u)$, dans l'ordre
inverse de la convention $(F_{0}, F_{1}, u)$ de la page 3. La justification
en une ligne est la nôtre : sur un recouvrement où $F$ est constant, $u$ est un
isomorphisme, donc partout.}

\emph{Un ouvert, c'est un morphisme vers le topos de Sierpiński.} Partons
inversement d'un topos $\mathcal{X}$ et d'un sous-topos ouvert
$\mathcal{X}_{0} \subset \mathcal{X}$, c'est-à-dire d'un sous-objet de l'objet
final de $\mathcal{F}(\mathcal{X})$, d'où le topos fermé résiduel
$\mathcal{X}_{1}$. La même donnée est celle d'un morphisme
\[
\pi : \mathcal{X} \longrightarrow \mathrm{Top}(\Delta_{1}),
\]
$\mathcal{X}_{0}$ étant l'image inverse du sous-topos ouvert « universel »
$\{0\}$ et $\mathcal{X}_{1}$ celle du fermé $\{1\}$.\footnote{Le topos de
Sierpiński classifie les sous-objets de l'objet final ; la formulation est sur
la page, le nom ne l'est pas.}

\emph{Recollements quelconques} (page 6). La donnée de
$(\mathcal{X}, \mathcal{X}_{0})$ équivaut à celle d'un triple
$(\mathcal{X}_{0}, \mathcal{X}_{1}, \varphi)$ où
$\varphi : \mathcal{F}(\mathcal{X}_{0}) \to \mathcal{F}(\mathcal{X}_{1})$ est un
foncteur exact à gauche quelconque. On cherche des conditions sur ce triple qui
assurent que les deux inclusions
\[
\mathcal{X}_{0} \longrightarrow \mathcal{X} \longleftarrow \mathcal{X}_{1}
\]
soient des équivalences faibles. Les cas qui intéressent le plus sont ceux où
$\varphi$ commute aux limites projectives quelconques, c'est-à-dire admet un
adjoint à gauche $f^{*}$ --- comme au paragraphe précédent, où l'on supposait de
plus $f^{*}$ exact à gauche, ce qu'on ne veut pas supposer ici.\footnote{La page
écrit ici « un adjoint à droite $f^{*}$ », là où les pages 4 et 7 écrivent « à
gauche » dans la même phrase. Un foncteur qui commute aux limites projectives
a un adjoint à gauche ; on corrige.} Le dossier ne donne pas ces conditions.

\emph{Le cas des préfaisceaux} (page 7). Supposons $\mathcal{X} =
\mathrm{Top}(X)$ pour une petite catégorie $X$. Un sous-objet de l'objet final
de $\widehat{X}$ est un crible, ensemble d'objets stable par les flèches qui y
arrivent ; il découpe $X$ en deux sous-catégories pleines $X_{0}$ (le crible)
et $X_{1}$ (le complémentaire), et $\mathcal{X}_{0} = \widehat{X_{0}}$,
$\mathcal{X}_{1} = \widehat{X_{1}}$. Le foncteur de recollement commute alors
aux limites quelconques, et son adjoint à gauche, qui commute aux limites
inductives, est déterminé par sa restriction
\[
X_{1} \longrightarrow \widehat{X_{0}},
\]
c'est-à-dire par un bifoncteur
\[
h : X_{0}^{\mathrm{op}} \times X_{1} \longrightarrow (\mathrm{Ens}),
\qquad (x_{0}, x_{1}) \longmapsto h(x_{1}, x_{0}),
\]
qui est $h(x_{1}, x_{0}) = \mathrm{Hom}_{X}(x_{0}, x_{1})$ ; il n'y a pas de
flèche de $X_{1}$ vers $X_{0}$, puisque $X_{0}$ est un crible.
Réciproquement, un bifoncteur $h$ quelconque définit une catégorie $X$
d'objets $\mathrm{Ob}\, X_{0} \sqcup \mathrm{Ob}\, X_{1}$, dont $\widehat{X}$ est
le recollement correspondant.\footnote{Le crible et la description de $h$ par
les flèches de $X$ sont les nôtres ; la page s'arrête sur le bifoncteur. Une
telle catégorie $X$ est ce qu'on appelle aujourd'hui le \emph{collage}, ou
cographe, du profoncteur $h$.}

\pagerange{9}{9}
\section*{Un feuillet séparé (page 9)}

Un carré de flèches $A' \to A$ au-dessus de $B' \to B$, et, pour $A$ filtrante
et $b$ un objet de $B$, une comparaison entre la catégorie $A_{/b}$ des objets
de $A$ au-dessus de $b$ et la fibre $A_{b}$ :
\[
\varprojlim_{A_{/b}} \bigl(F|_{A_{/b}}\bigr) \xrightarrow{\;\sim\;}
\varprojlim_{A_{b}} \bigl(F|_{A_{b}}\bigr)
\qquad \text{pour tout $F$ sur $A$ ou sur $A_{/b}$.}
\]
La page ne dit pas ce que doit être le foncteur $A \to B$. L'isomorphisme vaut,
par exemple, quand $A \to B$ est une catégorie précofibrée : l'inclusion de la
fibre $A_{b} \to A_{/b}$ admet alors un adjoint à gauche, donné par le
transport le long de la flèche structurale, et un foncteur qui admet un adjoint
à gauche ne change pas les limites projectives.\footnote{Cette hypothèse et
l'argument sont les nôtres. En marge, écrit verticalement : « $f$ dans Hot »,
le dernier mot étant lu avec doute.}

\end{document}
